\paragraph{A worked payment calculation.}
Take two fixed owners, $f(c)=\min\{c_1,c_2\}$, the ordered forms
$V=((1,0),(0,1))$, and $j=2$. Owner 2 reports $b_2=(0,2)$ throughout;
compare owner 1's bundle $(3,0)$ with its binding underreport $(1,0)$.
The report space remains the full two-owner bundle domain. At either displayed
profile, write $T=f(b_1+b_2)$. When owner 1 is excluded, the first form is
selected for either value of $B_2$, including the all-zero tie, so $p_1^q=T$.
When owner 2 is excluded, its capped contribution is $T$ exactly when owner 1
is present, so $p_2^q=(1-q)T$. Both endpoint marginals equal $(1-q)T$.
For example, at the full report and mask $(B_1,B_2)=(0,1)$, owner 1's
raw contribution $3$ is capped at $T=2$. Thus
\[
 r_1^q=T(1-q+q^2),\qquad r_2^q=T(1-q),\qquad
 (M_1,M_2)=\left(\frac{5T}{6},\frac{3T}{4}\right).
\]
\begin{center}
\begin{tabular}{cccccc}
\toprule
Owner 1's report & $T$ & $M_1$ & $M_2$ & Total payment & Subsidy\\
\midrule
$(3,0)$ & $2$ & $5/3$ & $3/2$ & $19/6$ & $7/6$\\
$(1,0)$ & $1$ & $5/6$ & $3/4$ & $19/12$ & $7/12$\\
\bottomrule
\end{tabular}
\end{center}
The underreport reduces owner 1's payment. Singleton production values are
zero and total payment is $19T/12\ge T$, so every coalition is covered at
both profiles. These payments illustrate the specified rule, not a necessary
subsidy: the efficient two-owner Shapley alternative discussed in
Proposition~\ref{prod:uncapped} pays $(1,1)$ and $(1/2,1/2)$ respectively.

